\cleardoublepage
\chapter*{Podsumowanie}
\addcontentsline{toc}{chapter}{Podsumowanie}

Celem niniejszej pracy było przeprowadzenie analizy wydajności i~efektywności wybranych algorytmów sortowania równoległego na procesorach wielordzeniowych. 
W~języku Python zaimplementowano Quick Sort, Merge Sort, Bucket Sort oraz Sample Sort w~wariantach sekwencyjnych i~równoległych. 
Do realizacji równoległości wykorzystano moduł \texttt{multiprocessing} oraz pamięć współdzieloną.
Implementacje poddano badaniom dla zróżnicowanych rozmiarów i~charakterystyk danych wejściowych oraz różnej liczby jednostek wykonawczych.

Przeprowadzone eksperymenty umożliwiły porównanie algorytmów pod względem czasu wykonania, 
wykorzystania CPU i~pamięci RAM, przyspieszenia oraz efektywności. 
Uzyskane wyniki wykazały, że obliczenia równoległe mogą skracać czas sortowania, 
jednak osiągane korzyści zależały od algorytmu i~liczby wykorzystywanych jednostek wykonawczych.
Nie zaobserwowano jednej konfiguracji, która byłaby najkorzystniejsza dla wszystkich badanych implementacji.

Spośród badanych algorytmów najkrótsze czasy wykonania osiągnał Bucket Sort, poprawiając rezultaty do 16 jednostek wykonawczych.
Najwyższe przyspieszenie uzyskano natomiast dla Sample Sort przy 8 jednostkach.
Quick Sort i~Merge Sort osiągały najlepsze wyniki przy niższym poziomie równoległości. 
Dalsze zwiększanie liczby jednostek nie zawsze skracało czas wykonania, a~jednocześnie odnotowywano spadek efektywności.

Na uzyskane rezultaty oddziaływały między innymi sposób podziału danych, wielkość zadań, 
zastosowane progi przejścia do wykonania sekwencyjnego oraz narzut związany z~tworzeniem i~synchronizacją procesów.
Większy rozmiar danych umożliwiał lepsze wykorzystanie równoległości, 
natomiast wpływ ich charakterystyki był różny dla poszczególnych algorytmów.

Wyniki przeprowadzonych badań pozwalają stwierdzić, że założony cel pracy został zrealizowany.
Uzyskane rezultaty pokazują, że samo zwiększanie liczby jednostek wykonawczych nie gwarantuje poprawy wydajności, 
a~efektywne wykorzystanie równoległości wymaga dopasowania jej poziomu zarówno do konkretnego algorytmu, jak i~do rozmiaru przetwarzanych danych.